\subsection{The three-division quartic and the CM character}

\begin{result}{E.49}{Proved}{The CM character vanishes on GT}
For the continuous character $\Delta$ of LM7 E.65, with its
specified signed frames, one has
\[
 \Delta|_{\widehat{\mathrm{GT}}}=0.
                                                               \tag{E.254}
\]
The element constructed in E.46 therefore does not belong to
$\widehat{\mathrm{GT}}$.
The comparison uses the actual three-division configuration,
the ordinary universal-curve and forgetting maps, and the
named branch-order comparison of E.44.  The finite computations
below certify explicit words and matrices; the proof of the
profinite comparison is given separately.
\end{result}
\begin{proof}
It suffices first to take a multiplier-one element
$F=(1,f)\in\widehat{\mathrm{GT}}$.  Write
$B=C_2(s)*C_3(r)$, $\tau=sr$, and
$F_2=\ker(B\to\mathfrak S_3)$, with generators
$x=\tau^{-2}$ and $y=sxs$.  The action and pure frames of
LM7 E.65 are
\[
 \alpha(\tau)=\tau,\quad
 \alpha(s)=\beta s\beta^{-1},\quad
 \alpha(r)=A r A^{-1},\qquad A,\beta\in\widehat F_2.
\]
Put $U=\left(\begin{smallmatrix}1&1\\0&1\end{smallmatrix}\right)$,
$L_0=\left(\begin{smallmatrix}1&0\\-1&1\end{smallmatrix}\right)$,
$S=UL_0U$, $R_+=UL_0$, and $R_0=-R_+$.
If $(\beta^{-1})^{\rm ab}=(\rho,\rho)$, the normalized matrices
are $V=U^{8\rho}$, $T=q_{12}(A)^{-1}V$, and
$W=q_{12}(\beta)^{-1}V$.  They are $I$ modulo four;
at three, $T$ centralizes $R_0$ and $W\in\langle S\rangle$.
The two bits are the sign $\epsilon_3(T)$ on
$\ker(R_0-I)$ and the exponent $\epsilon_2(W)$ modulo two
in $\langle S\rangle$.  Thus
$\delta(F)=\epsilon_3(T)+\epsilon_2(W)$.
We prove equality of the two bits by identifying the maps that
produce them.

\par\medskip\noindent\textbf{The intrinsic three-torsion orbit.}\quad
Let $H=(\mathbb Z/3)^2$ act freely by translations on an
oriented torus, and mark one regular orbit by its nine elements.
The identity label is called $0$.  In the full mixed mapping
class group let $N$ be the normalizer of this actual subgroup,
and let $N_0$ fix label $0$.  There are actual isomorphisms
\[
 1\longrightarrow H\longrightarrow N\longrightarrow
 \Gamma_{1,1}\longrightarrow1,
 \qquad
 N_0\xrightarrow[\text{forget eight labels}]{\ \sim\ }
 \Gamma_{1,1}.
                                                               \tag{E.255}
\]
To prove them, remove the nine marks and the one quotient mark.
The covering sequence is $F_{10}\triangleleft F_2$ with quotient
$H$, defined by abelianization modulo three.  Its kernel is
characteristic.  The center-free covering-extension argument of
E.44 identifies the covering $F_2$ with the inverse image of
$H$ in $\operatorname{Aut}(F_{10})$.
A normalizing mapping class therefore descends to an oriented
once-punctured torus mapping class, by punctured
Dehn--Nielsen--Baer.  Conversely every such downstairs mapping
class preserves this covering kernel and lifts.  The kernel of
descent is precisely $H$.  A unique deck translate of a lift
fixes label $0$, proving the assertion about $N_0$.
On linear torus representatives, quotienting by $H$ and
forgetting the other eight labels give the same integral
homology matrix.  Thus it is the stated ordinary forgetting
map that gives the second isomorphism.

Every finite quotient of $N_0$ consequently extends to the
appropriate finite-index ambient group by that forgetting map.
This proves its full induced topology; intersecting with the
label-$0$ stabilizer proves the same for $N$.
The finite-$p$ theorem of \cite[Theorem A]{ET_BZ2024}, here
with $p=3$, now identifies the completed normalizer with
$\widehat N$.  Its hypotheses hold by the genus-one goodness
and finite virtual cohomological-type arguments of D.35.

There is an intrinsic way to exclude a different $C_3^2$ action
containing rotations of a CM torus.  For the translation action,
the inverse image of $H$ in the universal-curve extension is
the torsion-free $\widehat F_2$.  A geometric action with an
unmarked fixed point has a nontrivial finite-order lift in this
extension.  All nine marked points form a free orbit, so any
fixed point of a nonidentity element is unmarked.  Preservation
of the actual universal-curve extension therefore preserves
the free action type.  Finite-subgroup realization
\cite{ET_Kerckhoff1983}, followed by the finite-$p$ theorem,
applies this test to the whole image of $H$.
The label-compatible free actions have two pure conjugacy
classes: their monodromy isomorphisms
$(\mathbb Z/3)^2\to H$ differ by the two cosets of
$\operatorname{SL}_2(\mathbb F_3)$ in
$\operatorname{GL}_2(\mathbb F_3)$.
A determinant-minus-one relabelling interchanges them and
fixes label $0$.  It can be represented by a braid supported
in a disk containing the other eight marks and avoiding $0$;
after those marks are forgotten its mapping class is the
identity.  This correction will only permute the four
target branch labels below, where it is a geometric mixed
braid.  Both classes, and this correction, are retained.

\par\medskip\noindent\textbf{Adding the fixed points and taking the actual quotient.}\quad
The center of $N_0\simeq\Gamma_{1,1}$ has the elliptic
involution $\iota$.  The nine original marks consist of its
fixed origin and four pairs.  The successive inverse images of
$\langle\iota\rangle$ when adding its three unmarked fixed
points are the completions of
\[
 (C_2*C_2*C_2)*F_4,\quad
 (C_2*C_2)*F_5,\quad C_2*F_6,\quad F_7.
                                                               \tag{E.256}
\]
Their even kernels are $F_{10},F_{11},F_{12},F_{13}$.
The reduced-word torsion theorem and the $p=2$ case of
Theorem A give respectively three, two, and one remaining
conjugacy classes of order-two lifts.  Each class under the
whole inverse-image group is one class under its even kernel:
its odd $C_2$ generator centralizes it.
The centralizer of a chosen lift intersects the free kernel
trivially.  Projection from that centralizer onto the class
stabilizer is onto, since an arbitrary lift can be corrected by
a fibre conjugator.  It is thus an isomorphism.  Its full
topology follows by pulling finite quotients back along
ordinary forgetting to the finite-index class stabilizer in
$N_0$.  The full Birman topology, or E.32's fibre-product
argument, identifies these with the actual completed
universal-curve pullbacks.
Successively taking the ordered groupoids of these lift
classes, then their finite permutation descent, therefore
adjoins all three points functorially.  This constructs the
twelve-point set $E[2]\cup E[3]$; no independent algebraic
torsion operation is assumed.

Here is the full topology check for the involution quotient
used next.  Let $G$ preserve its four fixed labels individually
and its four named pairs, allowing the independent pair flips.
Set $C=C_G(\iota)$.  Removing all twelve upstairs marks and all
eight downstairs marks gives the actual unramified cover
\[
 P_0=F_{13}\triangleleft E_0=F_7,\qquad E_0/P_0=C_2,
 \qquad
 1\longrightarrow\langle\iota\rangle\longrightarrow C
 \longrightarrow\Gamma_{0,8}\longrightarrow1.
                                                               \tag{E.257}
\]
The covering character is one on each branch cusp and zero
on each pair cusp.  To justify the second map, use the inverse
image of $\langle\iota\rangle$ in
$\operatorname{Aut}(P_0)$.  At a branch cusp, preservation of
the positive upstairs square and unique roots in $F_7$ imply
preservation of its positive downstairs root.  At a pair
cusp no square is needed.  Punctured Dehn--Nielsen--Baer
gives the pure downstairs mapping class; its kernel is the
deck group.  Conversely every pure downstairs mapping class
preserves the displayed parity and lifts.

Put $P=\widehat F_{13}$ and $E=\widehat F_7$.
Then $C_E(P)=1$: a centralizing element has square in
$Z(P)=1$, and $E$ is torsion-free.  Thus $E$ is exactly the
inverse image of $\langle\iota\rangle$ in
$\operatorname{Aut}(P)$, and restriction gives a continuous
map $C_{\operatorname{Out}(P)}(\iota)\to
\operatorname{Out}(E)$.
Completing (E.257) is right exact, and its finite kernel
survives because $\iota$ acts as $-I$ on torus homology
modulo three.  If an element of $\widehat C$ acts trivially
on $P$, its image in $\widehat\Gamma_{0,8}$ acts trivially
on $E$.  Genus-zero congruence faithfulness, proved in D.19,
puts it in the deck kernel, where it is again detected by
torus homology.  This proves that
$\widehat C\to\operatorname{Out}(P)$ is injective.
It factors through the actual ambient $\widehat G$, proving
the full induced topology of $C$.  The finite-$p$ centralizer
theorem therefore gives exactly
\[
 C_{\widehat G}(\iota)=\widehat C,
 \qquad
 C_{\widehat G}(\iota)/\langle\iota\rangle
       \xrightarrow{\ \sim\ }\widehat\Gamma_{0,8}.
                                                               \tag{E.258}
\]
Continuity of the restriction map follows on finite
characteristic quotients of the finitely generated free
groups; it is not an asserted intersection of closures.

There is one pure conjugacy class of involutions with these
fixed and paired labels.  The quotient is an eight-point
sphere; moving any one ordinary point around a branch point
flips its two lifts independently.  These four loops make
the cover of all individual sheet choices connected over
the connected $M_{0,8}$.  Finite-subgroup realization and
Theorem A therefore straighten the involution by a pure
conjugator.  This retains every named forgetting map.
Forgetting three pairs gives four actual maps
$\widehat\Gamma_{0,8}\to\widehat\Gamma_{0,5}$, each
retaining the four branch labels and its selected pair label.
By \cite[Corollary 2.8]{ET_HMM2022}, write the induced
eight-point outer automorphism as $(F',s)$ in
$\widehat{\mathrm{GT}}\times\mathfrak S_8$.
The four forgetting kernels force $s$ to preserve each
retained five-element set: a pair twist survives exactly
when both of its labels survive.  Hence it fixes all four
pair labels and only permutes the branch quartet.
E.44's pinned five-point comparison is $(F,1)$ on each
quotient.  One such quotient forces $F'=F$ and fixes the
branch quartet too.  Thus the eight-point parameter is
exactly $F$, with every label pinned.
Finite mixed-label descent is compatible as well: the
completed pure groups for eight and five points are
center-free and have trivial centralizer in their mixed
extensions, so an extension agreeing there is unique.

\par\medskip\noindent\textbf{The division quartic with its two ordered poles.}\quad
For $E:y^2=x^3+ax+b$, put $D=4a^3+27b^2\ne0$.
The four images of its nonzero three-torsion points are the
roots of $3x^4+6ax^2+12bx-a^2$.
Take the monic double cover and set
\[
\begin{gathered}
 t^2=x^4+2ax^2+4bx-a^2/3,
 \quad X=(t+x^2+a)/2,\quad Z=X-a/3,\quad Y=Xx+b/2,\\
 Y^2=Z^3+D/108,
 \qquad P=(-a/3,-b/2),\\
 x=\frac{Y-b/2}{Z+a/3},\qquad
 t=2(Z+a/3)-x^2-a.
\end{gathered}                                                   \tag{E.259}
\]
The two smooth projective curves are isomorphic by these
mutually inverse rational maps.  The positive pole is the
origin and the negative pole is $P$.  The quartic involution
is $Q\mapsto P-Q$, so its four branch points satisfy
$2Q=P$.  At $b=0$, $P$ is a nonzero two-torsion point;
at $a=0$ it is a nonzero three-torsion point, since its
horizontal tangent has triple intersection.
The original scaling by $u$ induces
$(Z,Y)\mapsto(u^4Z,u^6Y)$.  Thus $u=i$ induces the elliptic
involution, while $u=\zeta_6$ induces multiplication by
$\zeta_3^2$ in the displayed target coordinate.
These are explicit maps, including the order-three
orientation, and not a replacement based on the abstract
CM type.

The intrinsic construction just proved obtains the four
roots and the original origin as a five-point sphere:
take the twelve-point involution quotient and forget the
other three branch values.  Its ordered-pole quartic is
exactly E.44's named mixed five-point construction.
The ordering of its poles is an actual coordinate-independent
choice.  For an affine coordinate with distinguished point
$\infty$, use the monic polynomial
$t^2=\prod_i(x-\alpha_i)$ and define $p_\pm$ by
$t/x^2\to\pm1$.  Under $x'=cx+d$ the canonical map is
$t'=c^2t$, so the sign of $t'/x'^2$ is unchanged.
This also holds in families and around loops of affine
coordinates.  Every root relabelling preserves the monic
polynomial and both named poles.  Thus the lift is to the
group with its two elliptic labels individually fixed,
not to an unspecified extension allowing their exchange.
Thus (E.259) is compatible with the same parameter $F$.
The determinant-minus-one correction of the original
$H$-level basis only permutes these four roots.  It is
therefore an inner geometric braid in this exact
$\Gamma_{1,2}=B_4/Z(B_4)$ target: a root braid lifts by
the displayed monic construction and fixes $p_+,p_-$
individually.

More precisely, let $G_2=\widehat\Gamma_{1,2}$,
$G_1=\widehat\Gamma_{1,1}$ and let $q:G_2\to G_1$ be
ordinary forgetting.  Write $\mathcal E_0$ for the smooth
projective curve $Y^2=Z^3+1$, with origin $O$ at infinity.
The rigidified source and the target
CM locus are the isomorphic stacks
\[
 \left[\{(a,b):D\ne0\}/\mathbb G_m{}_{(2,3)}\right]
 \ \xrightarrow{\ \sim\ }\
 [(\mathcal E_0\setminus\{O\})/\mu_6],
 \qquad
 \varphi:\widehat B\xrightarrow{\ \sim\ }
 G_3:=q^{-1}(\langle R_+\rangle).
                                                               \tag{E.260}
\]
Indeed the level set $D=108$ is the indicated affine curve
under $(Z,Y)=(-a/3,-b/2)$, with
residual scaling group $\mu_6$.  Conversely take
$a=-3Z,b=-2Y$; then $D=108$.  This proves the stack
isomorphism and its fundamental-group map.
The source is rigidified because the original central
involution fixes its $x$-coordinate.  Explicitly, if $u$
is the original Weierstrass scaling, put $v=u^2$ after
rigidification.  Then $(a,b)$ have weights $(2,3)$,
$D$ has weight six, and $(Z,Y)$ transform by
$(v^2Z,v^3Y)$.  On the level $D=108$ this is precisely
the standard residual $\mu_6$ action, with the origin
and the second point still ordered.  The target is the
actual fibre over the CM stabilizer $\mu_6$ in the
universal curve, so its group is the displayed preimage.
The completion in (E.260) also has its full induced
topology.  The discrete CM pullback has exact sequence
$1\to F_2\to B\to C_6\to1$; the free kernel has its
full topology because it has finite index.  On the
ambient side, full Birman exactness embeds this same
$\widehat F_2$ as the pointpush kernel of
$G_2\to G_1$, and the finite CM subgroup $C_6$ embeds
in $G_1$.  The map from the completed pullback to
$q^{-1}(C_6)$ is the identity on this kernel and on
the finite quotient.  A kernel--quotient argument
therefore makes it injective and surjective.  This
proves the displayed completed isomorphism without
regarding the CM subgroup as finite-index in $G_2$.
In particular, $B'$ is identified with the pointpush
free group.  On the quotient, $\varphi(\tau)$ maps to
$R_+^{\varepsilon}$ for the actual
$\varepsilon\in\{1,-1\}$.
We retain that exponent below; no choice of positive
generator is changed without its map.

\par\medskip\noindent\textbf{The signed pointpush normalization.}\quad
Use E.11's filling
$B_4/\langle\omega_4\rangle\to
B_3/\langle\omega_3^2\rangle$ with
$(a,b,c)\mapsto(u,w,u)$.
Its free kernel $K$ has generators
$\xi_1=ca^{-1}$ and
$\xi_2=aba c^{-1}b^{-1}a^{-1}$, with exact conjugations
\[
\begin{array}{c|cc}
 &\xi_1&\xi_2\\ \hline
 a&\xi_1&\xi_2\xi_1\\
 b&\xi_1\xi_2^{-1}&\xi_2\\
 c&\xi_1&\xi_1\xi_2
\end{array}
 \qquad Q(a)=Q(c)=U,\quad Q(b)=L_0.
                                                               \tag{E.261}
\]
These relations also give the split presentation
$B_4=F(\xi_1,\xi_2)\rtimes B_3$, with $c=\xi_1a$.
The central quotient retains the free kernel because
$\omega_4$ maps to the infinite-order $\omega_3^2$
before either quotient.  Full completion is the Birman
sequence already identified in E.11.

Let $\Phi$ be the genuine four-strand action (E.13).
Put $e=f(ba^2b,c^2)$.  The identity
$ba^2b=(ab)^3a^{-2}$ and the commutation of $a$ with
both arguments give $e=f((ab)^3,c^2)$, by the
central-shift identity for a word in the commutator
subgroup.  Thus $[e,a]=1$.  Filling makes its two
arguments commute, so $e\in\widehat K$.
Since $\Phi(c)=e^{-1}ce$ and $\Phi(a)=a$, the matrix
$N$ induced on $K^{\rm ab}$ satisfies
\[
 (U-I)[e]=0,\qquad Ne_1=e_1,\qquad NU=UN,
 \qquad N=U^t,\qquad N\equiv V\pmod {12}.
                                                               \tag{E.262}
\]
The first three equations force the exact form $U^t$;
in particular they remove every possible scalar and
every primewise sign.  Covariance with $b$ gives
$NL_0N^{-1}=q_{12}(f)^{-1}L_0q_{12}(f)$.
The signed identity of LM7 E.65 gives the same equation
for $V$ modulo twelve.  It has also been recomputed here:
$H=\langle U^2,L_0^2\rangle$ has 96 elements and the
homomorphism $\epsilon:H\to C_3$ sends both generators
to one.  For $g=q_{12}(\beta)$ one has
$\epsilon(g)=-2\rho=\rho$ modulo three and
$z=q_{12}(f)=SgS^{-1}g^{-1}$.  The necessary conditions
$z=I$ modulo four and $(z^{-1}R_0)^3=I$ leave exactly
24 elements; all give
$z=\left(\begin{smallmatrix}1&-8\rho\\
-8\rho&1+64\rho^2\end{smallmatrix}\right)$.
The signed cube condition follows from its projective
version because reduction modulo four excludes $-I$.
The exhaustive Cayley enumeration and every matrix are
retained in \texttt{E\_CM\_covariance\_audit.py} and
its JSON certificate; it checks the displayed covariance
with its sign.  The common centralizer of
$U,L_0$ is scalar, and both matrices fix $e_1$, proving
the last equality in (E.262).

The canonical lift $\widetilde A$ of the pure frame to
$G_1$ is its word in $u^2,w^2$; its signed image is
$q_{12}(A)$.  One has
$\Phi(R_+)=\widetilde A R_+\widetilde A^{-1}$.
The equality modulo the central involution is the
defining frame equation.  The other possible sign is
excluded by $G_1^{\rm ab}=\mathbb Z/12$, on which the
multiplier-one action is the identity.
Lift the same word to $G_2$ using $a^2,b^2$.
The corrected automorphism
\[
 \Phi^A=\operatorname{Ad}(\widetilde A^{-1})\Phi
 \quad\text{preserves }G_3,\qquad
 (\Phi^A)_{K^{\rm ab}}\equiv
          q_{12}(A)^{-1}V=T\pmod {12}.
                                                               \tag{E.263}
\]
This proves the named normalization of $T$ directly
from the pointpush formulas.

\par\medskip\noindent\textbf{The commutator homology of the modular source.}\quad
Put $K_0=B'$.  Its abelianized extension is explicitly
\[
 B/[K_0,K_0]=\mathbb Z^2\rtimes C_6,\qquad
 R=\begin{pmatrix}0&-1\\1&1\end{pmatrix},\quad
 \tau=(0,R),\quad s=(p,R^3),\quad r=(p,R^4),\quad p=e_1.
                                                               \tag{E.264}
\]
Here is a direct Schreier verification.  In $K_0^{\rm ab}$
put $p_i=\tau^i(s\tau^{-3})\tau^{-i}$ and
$z=\tau^6$.  The two torsion relations give
$p_i+p_{i+3}+z=0$ and
$p_i+p_{i+4}+p_{i+8}+2z=0$, with indices periodic
modulo six.  The equations at $i=0,1$ force $z=0$
and $p_2=p_1-p_0$.  The others then follow.
The displayed semidirect product shows $p_0,p_1$
independent, with conjugation matrix $R$.
The kernel $K_0$ is a torsion-free index-six subgroup
of $C_2*C_3$, hence free of rank two.  Finite-index
completion gives the same calculation over
$\widehat{\mathbb Z}$.

Let $M$ be the matrix of $\alpha$ on $K_0^{\rm ab}$.
It commutes with $R$, since $\alpha(\tau)=\tau$.
Write the two pure frames in (E.264) as
$\beta=(b,R^j)$ and $A=(a,R^\ell)$.
Both $j$ and $\ell$ are even, because their images
in the abelianization of $\mathfrak S_3$ are zero.
The exact frame equations are
\[
 Mp=2b+R^jp,
 \qquad Mp=(I-R^4)a+R^\ell p,
 \qquad j=4\rho\quad\text{in }\mathbb Z/6.
                                                               \tag{E.265}
\]
The last equality follows from the exponent $-2$ of
each of $x=\tau^{-2}$ and $y=sxs$ and the exponent
vector $(-\rho,-\rho)$ of $\beta$.

Modulo three, $R^4=-R$ and $R+I$ is nonzero nilpotent
of rank one, whose image and kernel are the same
line $L_-$.  The second equation of (E.265) shows
that $M$ acts as $+1$ on the quotient by this line,
since $\ell$ is even.  Commutation with $R+I$ then
gives the same scalar on the line itself.
To compute at four, identify the projective action
$q_3:B\to\operatorname{PSL}_2(\mathbb F_3)=A_4$
with the affine action on $\mathbb F_2^2$ by
\[
 \infty\mapsto0,\quad 0\mapsto p,\quad
 -1\mapsto Rp,\quad 1\mapsto p+Rp.
\]
The named elements $S,R_0,U^{-1}$ act as
$s=(p,I)$, $r=(p,R^4)$, and $\tau=(0,R)$.
By $j=4\rho$, the affine image of $V_3=U^{8\rho}$
is $(0,R^j)$.  Thus the projective image of
$W=q_3(\beta)^{-1}V_3$ is translation by
$-R^{-j}b$.  Since $W\in\langle S\rangle$, it is
translation by $\epsilon_2(W)p$.
Consequently $b=\epsilon_2(W)R^jp$ modulo two.
Using (E.265) and then commutation with $R$ on the
basis $p,Rp$ proves
\[
 M|_{L_-}=1\pmod3,
 \qquad M\equiv(-1)^{\epsilon_2(W)}R^j\pmod4,
 \qquad j\equiv0\pmod2.
                                                               \tag{E.266}
\]
In particular, the exponent of $M$ in the order-six
group $\langle R\bmod4\rangle$ has parity
$\epsilon_2(W)$.

\par\medskip\noindent\textbf{The single global CM gauge.}\quad
Naturality of the constructed map (E.260) supplies
an element $h\in G_2$ such that
\[
 \Phi^A\varphi=\operatorname{Ad}(h)\varphi\alpha,
 \qquad q(h)=R_+^k\quad(k\in\mathbb Z/6).
                                                               \tag{E.267}
\]
For the second assertion, both sides of the first
equation induce the identity on $\langle R_+\rangle$.
Therefore $q(h)$ centralizes $R_+$ in $G_1$.
Modulo its central involution, this is the centralizer
of the finite $C_3$ factor in $\widehat C_2\amalg
\widehat C_3$.  The $p=3$ theorem gives exactly
$C_3$, whose preimage is $\langle R_+\rangle$.
Hence $h\in G_3$ and there is one global exponent
$k$, rather than independent signs at different primes.

Let $C$ denote the integral homology isomorphism
induced by the actual $\varphi$.
It sends $R$ to $R_+^{\varepsilon}$, so sends the
line $L_-$ at three to $\ker(R_0-I)$.
Equations (E.263) and (E.267) give
\[
 TC=R_+^kCM\pmod{12},\qquad
 \epsilon_3(T)\equiv k\equiv\epsilon_2(W)\pmod2.
                                                               \tag{E.268}
\]
Indeed $M$ is $+1$ on the line at three, while
$R_+$ acts there by $-1$; this proves the first
parity equality.  At four $T=I$, so the CM exponent
of $M$ has parity $k$.  Equation (E.266) proves the
second.  The actual $\varepsilon=\pm1$ affects no
parity, and has remained in the comparison throughout.
Thus $\delta(F)=0$ for every multiplier-one GT element.

For arbitrary $F\in\widehat{\mathrm{GT}}$, choose a
Galois element $\sigma$ of the same cyclotomic
multiplier.  The element $F\sigma^{-1}$ is again in
GT and has multiplier one.  The defining, choice-independent
extension formula of LM7 E.65 now gives
\[
 \Delta(F)=\delta(F\sigma^{-1})=0,
 \qquad \Delta(F_{\mathrm{E.46}})=1.
                                                               \tag{E.269}
\]
This proves both assertions of (E.254).
For the particular witness, the obstruction is already visible
in the named homological diagram: $\rho=0$, $T=5I$ and
$W=I$ give $M=I$ modulo four and $M|_{L_-}=1$ modulo
three.  Equation (E.268) would force $k$ even at four
and odd at three.  Thus no global CM gauge exists even
for this finite homological shadow.  Inertia preservation,
$\mathfrak S_3$ compatibility, multiplier one and trivial
peripheral homology remain the passing tests of E.24;
the pentagon and the ordinary-tower extension are now
excluded by the proved obstruction.  No sufficiency
of $\Delta=0$ for GT membership is asserted.
The programs \texttt{b\_cm\_quartic.py} and
\texttt{b\_cm\_module.py} verify the mutually inverse
quartic equations, the CM weights, all six exact
pointpush braid words, the named projective-affine
dictionary, all 24 modulo-four frame cases, and all
27 modulo-three frame cases.  The subgroup completions,
parameter comparison, and global gauge are supplied
by the proofs above, not by those finite checks.
\end{proof}
